\section{Ownership Contract}

The RCTX boundary spec records the dependency rule: RPLAN, RCOUNT, RHIST, and
RMAP may depend on RCTX concepts, but RCTX must not depend on those domain
packages \citep{rctx-boundary-2026}. The civic evidence package family spec
uses the same split: RCTX is the canonical machine context, and RMAP is the
public map presentation layer.

\begin{center}
\small
\begin{tabularx}{\textwidth}{lX}
\toprule
Owner & Contract \\
\midrule
RCTX & Stable unit ids, unit namespace, canonical order, context hash, unit-universe hash, graph identity, source hashes, aligned non-domain attributes, and crosswalks. \\
RMAP & Rendered maps, styling, symbols, colors, labels, layer composition, display projections, raster/vector tiles, and public map image hashes. \\
RPLAN & District assignments, plan metadata, plan audit certificates, and the current \texttt{.rctx} read/write path until another full-context producer needs extraction. \\
RCOUNT & Election totals, CVRs, canvass records, ballot manifests, audit evidence, and references to RCTX context or crosswalk hashes through count-package metadata. \\
RHIST & Unit lineage across cycles: splits, merges, renames, boundary edits, and historical cross-cycle relationships. \\
\bottomrule
\end{tabularx}
\end{center}

The implemented \texttt{rctx-core} record shape follows that ownership line.
\texttt{RctxPackage} contains a \texttt{RctxManifest}, a source index, context
unit indexes, graph records, crosswalk records, and a claim boundary.
\texttt{RctxManifest} names the RCTX version, package id, jurisdiction,
producer, creation time, and \texttt{package\_content\_hash}. A
\texttt{ContextUnitIndex} binds a \texttt{context\_hash} to ordered
\texttt{unit\_ids}. A \texttt{GraphRecord} binds graph identity to a context
hash, graph hash, and source references. A \texttt{CrosswalkRecord} binds
source and target contexts, endpoint units, an exact rational weight, a weight
kind, an exhaustiveness flag, and source references.

The existing RPLAN context maps directly onto the shared boundary:

\begin{center}
\small
\begin{tabularx}{\textwidth}{lX}
\toprule
RCTX concept & Existing \texttt{rplan\_core::RplanContext} field \\
\midrule
Schema version & \texttt{rctx\_version} \\
Context digest & \texttt{context\_hash} \\
Canonical unit universe & \texttt{units: PlanUnitIndex} \\
Unit kind, state, year, order, source id & \texttt{PlanUnitIndex} fields \\
Unit-universe hash & \texttt{unit\_universe\_hash} \\
Adjacency graph & \texttt{graph: Option<UnitGraph>} \\
Aligned population vector & \texttt{populations: Option<Vec<i64>>} \\
Aligned subdivision ids & \texttt{subdivisions: Option<SubdivisionContext>} \\
Aligned demographic vectors & \texttt{demographics: Option<DemographicContext>} \\
Geometry identity references & \texttt{geometry: Option<GeometryContext>} \\
Source registry & \texttt{source\_hashes: SourceHashes} \\
\bottomrule
\end{tabularx}
\end{center}

That mapping is why X.2 does not move the full \texttt{.rctx} implementation.
The as-built RPLAN context already carries the fields needed by the shared
contract; the extraction point is the cross-product verification primitive, not
a rewrite of working plan context I/O.
